% =====================================================================
%  preamble.tex -- packages, styles, colours and macros
%  Kept to standard TeX Live / Overleaf packages so the document builds
%  first-try with pdflatex (no shell-escape, no minted).
% =====================================================================

% ---- Encoding & fonts ------------------------------------------------
\usepackage[T1]{fontenc}
\usepackage[utf8]{inputenc}
\usepackage{lmodern}
\usepackage{microtype}

% ---- Page geometry & headers ----------------------------------------
\usepackage[a4paper,margin=1in,headheight=14pt]{geometry}
\usepackage{fancyhdr}
\usepackage{lastpage}

% ---- Maths & symbols -------------------------------------------------
\usepackage{amsmath}
\usepackage{amssymb}
\usepackage{siunitx}
\sisetup{detect-weight=true,detect-family=true}

% ---- Tables ----------------------------------------------------------
\usepackage{booktabs}
\usepackage{array}
\usepackage{multirow}
\usepackage{tabularx}
\usepackage{longtable}
\usepackage{colortbl}
\usepackage{makecell}

% ---- Graphics & figures ---------------------------------------------
\usepackage{graphicx}
\graphicspath{{figures/}}
\usepackage{float}
\usepackage{caption}
\usepackage{subcaption}
\captionsetup{font=small,labelfont=bf}
\captionsetup[sub]{font=footnotesize}

% ---- Colour ----------------------------------------------------------
\usepackage[dvipsnames,table]{xcolor}
\definecolor{accent}{HTML}{1F4E79}      % deep blue -- titles / rules
\definecolor{accentlight}{HTML}{DAE3F3} % light blue -- table headers
\definecolor{arithblk}{HTML}{2E75B6}    % arithmetic datapath
\definecolor{logicblk}{HTML}{548235}    % logic datapath
\definecolor{muxblk}{HTML}{BF8F00}      % multiplexers
\definecolor{padblk}{HTML}{C55A11}      % I/O pads
\definecolor{coreblk}{HTML}{7030A0}     % core
\definecolor{codebg}{HTML}{F6F8FA}
\definecolor{codekw}{HTML}{0000C0}
\definecolor{codecomment}{HTML}{2E8B57}
\definecolor{codestr}{HTML}{A31515}
\definecolor{codenum}{HTML}{B36B00}

% ---- Drawing (TikZ / CircuiTikz) ------------------------------------
\usepackage{tikz}
\usetikzlibrary{arrows.meta,positioning,calc,fit,backgrounds,
                shapes.geometric,shapes.misc,decorations.pathreplacing,
                decorations.markings,chains}
\usepackage[siunitx,RPvoltages]{circuitikz}
\ctikzset{logic ports=ieee}

% Reusable TikZ styles for block diagrams -----------------------------
\tikzset{
  blk/.style     = {draw,thick,rounded corners=2pt,minimum height=9mm,
                     minimum width=16mm,align=center,font=\small},
  arith/.style   = {blk,draw=arithblk,fill=arithblk!12,text=black},
  logic/.style   = {blk,draw=logicblk,fill=logicblk!12,text=black},
  muxb/.style    = {draw=muxblk,fill=muxblk!15,thick,align=center,font=\small},
  core/.style    = {blk,draw=coreblk,fill=coreblk!10,minimum width=26mm,
                     minimum height=18mm},
  pad/.style     = {draw=padblk,fill=padblk!18,thick,minimum width=4mm,
                     minimum height=4mm,inner sep=0pt},
  sig/.style     = {-{Latex[length=2mm]},thick},
  bus/.style     = {-{Latex[length=2.4mm]},line width=1.1pt},
  buslbl/.style  = {font=\scriptsize\ttfamily,inner sep=1pt,fill=white},
  flow/.style    = {blk,minimum width=42mm,minimum height=11mm,
                     draw=accent,fill=accentlight,text=black},
  flowarr/.style = {-{Latex[length=2.6mm]},line width=1pt,draw=accent},
}

% Trapezoidal 4:1 multiplexer symbol: \muxfour{<name>}{<x>}{<y>}
% Draws a downward-tapering trapezoid with four inputs (0..3) on the left.
\newcommand{\muxfour}[3]{%
  \begin{scope}[shift={(#2,#3)}]
    \fill[muxblk!15] (0,0) -- (1.1,-0.35) -- (1.1,-1.85) -- (0,-2.2) -- cycle;
    \draw[muxblk,thick] (0,0) -- (1.1,-0.35) -- (1.1,-1.85) -- (0,-2.2) -- cycle;
    \node[font=\scriptsize] at (0.55,-1.1) {\shortstack{4:1\\MUX}};
    \coordinate (#1-in0) at (0,-0.30);
    \coordinate (#1-in1) at (0,-0.80);
    \coordinate (#1-in2) at (0,-1.30);
    \coordinate (#1-in3) at (0,-1.80);
    \coordinate (#1-out) at (1.1,-1.10);
    \coordinate (#1-sel) at (0.55,-2.03);
  \end{scope}%
}

% ---- Source code listings (Verilog) ---------------------------------
\usepackage{listings}
\lstdefinelanguage{verilogx}[]{Verilog}{%
  morekeywords={wire,reg,logic,always,assign,module,endmodule,input,output,
                inout,begin,end,case,casez,casex,endcase,default,parameter,
                localparam,generate,endgenerate,genvar,for,if,else,posedge,
                negedge}}
\lstdefinestyle{vlog}{
  language=verilogx,
  backgroundcolor=\color{codebg},
  basicstyle=\ttfamily\footnotesize,
  keywordstyle=\color{codekw}\bfseries,
  commentstyle=\color{codecomment}\itshape,
  stringstyle=\color{codestr},
  numberstyle=\tiny\color{gray},
  numbers=left,numbersep=7pt,
  showstringspaces=false,
  breaklines=true,
  frame=single,rulecolor=\color{gray!40},
  framexleftmargin=6pt,
  tabsize=4,
  columns=fullflexible,
  keepspaces=true,
  captionpos=b,
  aboveskip=8pt,belowskip=6pt,
}
\lstset{style=vlog}

% ---- Hyperlinks (load late) -----------------------------------------
\usepackage{hyperref}
\hypersetup{
  colorlinks=true,
  linkcolor=accent,
  citecolor=accent,
  urlcolor=accent,
  filecolor=accent,
  pdftitle={Design of a 32-bit ALU: RTL-to-GDS ASIC Flow},
  pdfauthor={Salwan Aldhahab},
}
\usepackage[nameinlink]{cleveref}
% Give cleveref a name for the listings float (\lstinputlisting labels).
\crefname{lstlisting}{Listing}{Listings}
\Crefname{lstlisting}{Listing}{Listings}
\usepackage{enumitem}
\setlist{nosep,leftmargin=1.5em}

% ---- Header / footer style ------------------------------------------
\pagestyle{fancy}
\fancyhf{}
\renewcommand{\headrulewidth}{0.4pt}
\renewcommand{\footrulewidth}{0pt}
\fancyhead[L]{\small\itshape 32-bit ALU RTL-to-GDS}
\fancyhead[R]{\small\leftmark}
\fancyfoot[C]{\small\thepage}

% Chapter-page style
\fancypagestyle{plain}{%
  \fancyhf{}\renewcommand{\headrulewidth}{0pt}%
  \fancyfoot[C]{\small\thepage}}

% ---- Convenience macros ---------------------------------------------
% Screenshot from the term-project report:
%   \reportfig{<width frac>}{<file>}{<caption>}{<label>}
\newcommand{\reportfig}[4]{%
  \begin{figure}[htbp]\centering
    \includegraphics[width=#1\linewidth]{#2}%
    \caption{#3}\label{#4}%
  \end{figure}}

% Same, with a trailing "(Source: term-project report, Fig. N)" note.
\newcommand{\reportfignote}[5]{%
  \begin{figure}[htbp]\centering
    \includegraphics[width=#1\linewidth]{#2}%
    \caption[#3]{#3\\{\footnotesize\itshape #5}}\label{#4}%
  \end{figure}}

\newcommand{\code}[1]{\texttt{\small #1}}
\newcommand{\sel}{\code{S3\,S2\,S1\,S0}}
\newcommand{\um}{\si{\micro\meter}}
\newcommand{\umsq}{\si{\micro\meter\squared}}

% Signal / bus label helper for prose
\newcommand{\bus}[1]{\texttt{#1}}
